% ============================================================
% RSC Advances Conclusions — honest benchmark (repositioned 2026-06-23)
% ============================================================

\section{Conclusions}\label{sec:conclusions}

We measured how far cheap structure-based machine learning reaches against the
experimental singlet--triplet gap $\Delta E_\text{ST}$ of TADF emitters, with every
number reproducible from the deposited code. The study establishes where the field's
data, not its models, set the limit on predicting this quantity, and it isolates two
effects that follow from that limit. Three findings stand out.

\paragraph{Data quality, not model capacity, bounds accuracy.}
Random forests built from the 2D structure alone---Morgan fingerprints
(MAE = \qty{0.091}{\electronvolt}) and RDKit descriptors
(\qty{0.100}{\electronvolt})---predict experimental $\Delta E_\text{ST}$ as well
as NTO descriptors derived from a semi-empirical excited-state calculation
(\qty{0.096}{\electronvolt}), over \num{231} molecules and \num{212} scaffolds, and
the result holds across scaffold, cluster and random cross-validation. The two
representations agree to within \qty{0.017}{\electronvolt} on paired folds, so the
excited-state step buys no accuracy. The representations separate, however, on
ranking under laboratory transfer: under a DOI-grouped split the NTO model retains
Spearman $\rho = \num{0.35}$ while Morgan drops to \num{0.16}
($\Delta\rho = \num{0.19}$, \qty{95}{\percent} CI \numrange{0.04}{0.34}, excluding zero
only narrowly and on a single split). The excited-state features do not improve the
value of the gap; they preserve the \emph{order} of candidates across source
laboratories---a dissociation between accuracy and ranking that, to our knowledge, has
not been reported for this property, and one that a wider group-level interval could
still absorb.
SHAP on the NTO variant ties the signal to hole--electron overlap and
charge-transfer character, keeping the physics-informed representation interpretable.

The label bounds what any predictor can reach: independent reports of the same compound
scatter by \qty{0.072}{\electronvolt} RMS, and the per-molecule median actually
regressed on is known only to \qty{0.049}{\electronvolt}---about half the model error,
so a real term in the budget but not the whole of it. Quantified across \num{231}
molecules and \num{173} source papers, this is, to our knowledge, the first
inter-laboratory variability budget for the TADF singlet--triplet gap. Error is lower on molecules whose replicate reports agree
(\qty{0.074}{\electronvolt} against \qty{0.106}{\electronvolt}), but that ordering is
reproduced by a constant predictor and reflects the spread of the target within each
stratum rather than label quality, so we draw no conclusion from it. At this data scale, capacity is not the
missing ingredient: no higher-capacity learner we tested beats the random forest---a
frozen pretrained molecular-transformer embedding reaches only MAE
\qtyrange{0.109}{0.113}{\electronvolt}---and a \num{609}-core-hour range-separated
TD-DFT reference gap does not either. Separately,
the computed reference is itself functional-dependent (mean \qty{0.30}{\electronvolt},
up to \qty{0.58}{\electronvolt} between B3LYP and CAM-B3LYP); that bounds the meaning of
a \emph{computed} gap rather than the error of a model fitted to measured labels, and we
do not combine the two into one floor. Claims of near-DFT precision against such references
are not physically meaningful.

The leakage demonstration is not TADF-specific. Regressing $\Delta E_\text{ST}$ onto a
feature set that contains the $S_1$ and $T_1$ excitation energies reproduces the target
by construction ($\Delta E_\text{ST}\equiv E_{S_1}-E_{T_1}$); under cross-validation a
linear model given those energies recovers the target exactly ($R^2 = 1.000$, MAE $= 0$)
against $R^2 = 0.36$ without them, so the apparent accuracy of such models is arithmetic,
not learning. Any target that is an algebraic combination of features already in the
descriptor set will be ``predicted'' by arithmetic, invisibly to cross-validation. For a
field in which many reported ML results target a computed gap, this is a reminder to
benchmark against the measured quantity with features that exclude the constituent
energies.

From the 2D structure alone, with no quantum chemistry, the model acts as a coarse
triage filter---enriching low-gap candidates \numrange{1.2}{1.4}$\times$ over random
selection on held-out labelled data---useful for prioritisation but far from a
selector, and not a quantitative photophysical method. Where ranking must survive a
change of source laboratory the semi-empirical descriptors are the better choice,
the one place in this study where the excited-state step pays for itself
(\cref{subsec:benchmark}). On the filtered library
the conformal prediction intervals are too wide for molecule-level selection, so
no ranked shortlist is deposited. Matching a specific lighting niche (horticultural or
human-centric) needs the emission wavelength and oscillator strength, which this
scalar-gap model does not predict, and is left to downstream work.
We make no device-level performance claims: the spin--orbit-coupling-based
$k_\text{RISC}$ and external-quantum-efficiency estimates explored during this work
were molecule-independent and have been withdrawn.

\paragraph{What would push the field past the ceiling.}
The \qty{0.049}{\electronvolt} precision of the median label is the term a better
dataset, not a better model, would attack. A curated, single-source,
solvent-controlled experimental benchmark---one host, one measurement protocol, one
laboratory per compound, along the lines the photophysics community is already
converging on for how $\Delta E_\text{ST}$ should be reported~\cite{Shuaib2026}---would
reduce the inter-report scatter that currently \emph{is} the ceiling, and would let a
retrained model and a stronger representation be distinguished from label noise. We do not run that benchmark here; we provide the
quantitative evidence that it is the necessary next step. All data, trained models,
and the scripts that regenerate every number and figure in this work are openly
available (Zenodo DOI: 10.5281/zenodo.17436069), so that this benchmark can be
extended and the limitations we report can be tested directly.
